\documentclass[runningheads]{llncs}

% ---------------------------------------------------------------
% Include basic ECCV package
 
% TODO REVIEW: Insert your submission number below by replacing '*****'
% TODO FINAL: Comment out the following line for the camera-ready version
\usepackage[review,year=2026,ID=9159]{eccv}
% TODO FINAL: Un-comment the following line for the camera-ready version
%\usepackage{eccv}

% OPTIONAL: Un-comment the following line for a version which is easier to read
% on small portrait-orientation screens (e.g., mobile phones, or beside other windows)
%\usepackage[mobile]{eccv}


% ---------------------------------------------------------------
% Other packages

% Commonly used abbreviations (\eg, \ie, \etc, \cf, \etal, etc.)
\usepackage{eccvabbrv}

% Include other packages here, before hyperref.
\usepackage{graphicx}
\usepackage{booktabs}
\usepackage{multirow}

% The "axessiblity" package can be found at: https://ctan.org/pkg/axessibility?lang=en
\usepackage[accsupp]{axessibility}  % Improves PDF readability for those with disabilities.

% ---------------------------------------------------------------
% Hyperref package

% It is strongly recommended to use hyperref, especially for the review version.
% Please disable hyperref *only* if you encounter grave issues.
% hyperref with option pagebackref eases the reviewers' job, but should be disabled for the final version.
%
% If you comment hyperref and then uncomment it, you should delete
% main.aux before re-running LaTeX.
% (Or just hit 'q' on the first LaTeX run, let it finish, and you
%  should be clear).

% TODO FINAL: Comment out the following line for the camera-ready version
% \usepackage[pagebackref,breaklinks,colorlinks,citecolor=eccvblue]{hyperref}
% TODO FINAL: Un-comment the following line for the camera-ready version
\usepackage{hyperref}

% Support for ORCID icon
\usepackage{orcidlink}

% Graphics path
\graphicspath{{images/}}

\begin{document}

% ---------------------------------------------------------------
\title{EIHMR: Collaborative Human-Camera Estimation for Global Human Mesh Recovery}

\titlerunning{EIHMR}

% TODO FINAL: Replace with your author list.
\author{First Author\inst{1} \and
Second Author\inst{1} \and
Third Author\inst{1}}

\authorrunning{F.~Author et al.}

% TODO FINAL: Replace with your institution list.
\institute{Institution Name, City, Country\\
\email{email@example.com}}

\maketitle


% ===================================================================
% ABSTRACT
% ===================================================================
\begin{abstract}
Recovering global 3D human motion from monocular video captured by a moving camera is a fundamental yet challenging problem, as camera ego-motion and human body motion are tightly entangled in the image observations. The prevailing two-stage paradigm treats camera estimation and motion reconstruction as isolated processes, causing errors on both sides to be further amplified when combined in the world coordinate system. To address this, we draw inspiration from the human inner visual simulation mechanism and propose EIHMR, a collaborative human-camera co-estimation framework. EIHMR comprises two complementary stages that form a decouple-then-collaborate loop: Scene-Aware Local Human Motion Reconstruction reprojects motion sequences into a frozen keyframe viewpoint and applies metric depth and kinematic constraints to produce geometrically consistent local motion; Motion-aware SLAM then re-renders the refined motion as static meshes into the original frames, converting dynamic human regions into structured matching cues that strengthen camera estimation. On the EMDB dataset, EIHMR reduces global trajectory error by approximately 10\%, demonstrating strong capability in long-range human motion reconstruction.
\keywords{Human motion recovery \and Visual SLAM \and 3D human pose estimation}
\end{abstract}


% ===================================================================
% 1. INTRODUCTION
% ===================================================================
\section{Introduction}
\label{sec:intro}
Recovering global 3D human motion from monocular video constitutes a fundamental problem in computer vision. Beyond estimating kinematic body poses, this task requires recovering complete human trajectories within a world coordinate system to support downstream applications such as human--scene interaction understanding~\cite{Hassan2019,RICH2022}, embodied intelligence~\cite{Embodied2022,EmbodiedScan2024}, and digital human animation~\cite{DigitalHuman2024}. The challenge is especially acute when the input is captured by a moving camera, as camera motion and human body motion are tightly entangled in the image observations, rendering their separation inherently ambiguous.

The prevailing paradigm addresses this challenge through a two-stage pipeline: camera trajectories are first estimated via visual SLAM, after which human body poses are independently reconstructed and transformed into world coordinates~\cite{TRAM2024,SLAHMR2023,PACE2024}. However, this paradigm treats camera estimation and motion reconstruction as \textit{isolated} processes, and this isolation introduces two critical limitations. On the camera side, SLAM must mask out dynamic human regions to satisfy the static-scene assumption~\cite{DynaSLAM2018}; when the subject dominates the field of view, the available visual features are drastically reduced, severely degrading camera estimation quality. On the motion side, body poses estimated without access to scene geometry lack physically grounded spatial constraints, resulting in non-physical artifacts such as foot sliding, ground penetration, and floating limbs. These two sources of error are further amplified when combined in the world coordinate system.

We posit that overcoming these limitations necessitates a collaborative human--camera estimation paradigm. Our approach draws inspiration from the human cognitive process of perceiving others' motion. When observing a moving person, humans do not passively register visual impressions; rather, they mentally anchor a stable viewpoint, leverage scene cues such as foot--ground contact to reconstruct a scale-consistent local motion sequence, and subsequently employ their kinematic understanding of the human body to localize themselves and reconstruct the surrounding environment~\cite{shivkumar2025hierarchical,giese2003neural}. This \textit{decouple-then-collaborate} cognitive loop motivates the design of our framework: camera information provides scene anchors for the estimated motion sequences, facilitating scene-aware motion refinement, and the refined motion, in turn, enhances camera localization. However, achieving this is non-trivial, as it necessitates resolving the circular dependency between reconstructing scene-consistent motion from noisy camera estimates and leveraging that refined motion to stabilize SLAM against reciprocal error propagation.
% 缺少一句话说明难点，请修改下面llm生成的语句。However, non-trivial: local motion recovery depends on depth scale from the camera, while robust camera localization in human-dominated scenes requires geometrically consistent meshes that are themselves the output of motion recovery. Naively alternating the two stages propagates errors in both directions, demanding carefully designed interfaces to stabilize this mutual dependency.

To address these challenges, we propose EIHMR (\textbf{E}mbodied \textbf{I}magination \textbf{H}uman \textbf{M}esh \textbf{R}ecovery), a collaborative framework that emulates the human ability to mentally simulate observed motion in imagination. EIHMR comprises two complementary stages: Scene-Aware Local Human Motion Reconstruction (SA-LHMR) and Motion-aware SLAM (MO-SLAM). SA-LHMR reprojects motion sequences into a frozen keyframe viewpoint, thereby eliminating camera ego-motion interference, and applies metric depth and kinematic constraints via sliding-window optimization to produce geometrically consistent local motion free of non-physical artifacts. MO-SLAM re-renders the refined motion as geometrically consistent static meshes and injects them into the original frames, converting the human body from an excluded dynamic region into a source of dense, kinematically-aware matching cues that substantially strengthen camera estimation, particularly in human-dominated and texture-scarce scenarios. Together, the framework realizes a decouple-then-collaborate loop that supersedes the isolated, unidirectional information flow of prior work~\cite{WHAM2024,TRACE2023,GLAMR2022,GVHMR2024}. Our contributions are summarized as follows:

\begin{enumerate}
\item We introduce EIHMR, a collaborative human--camera co-estimation framework that first decouples human and camera motion in a frozen-camera reference frame for scene-aware refinement, and then feeds the refined motion back to augment camera estimation, establishing a bidirectional information flow absent in existing methods.

\item We propose SA-LHMR, which reprojects motion into a unified keyframe viewpoint and applies a two-stage refinement: root trajectory correction via forward-kinematics stationarity analysis, followed by contact-aware inverse kinematics guided by metric depth, effectively eliminating foot sliding and scene penetration artifacts.

\item We propose MO-SLAM, which re-renders refined human motion as geometrically consistent static meshes and integrates them into the SLAM pipeline, transforming the human body from a source of interference into structured features that provide dense matching cues and markedly improve camera estimation robustness.

\item We validate EIHMR through comprehensive experiments, demonstrating that the proposed collaborative strategy yields consistent gains across three complementary evaluation dimensions: camera-coordinate body pose estimation, camera trajectory recovery in dynamic human-dominated scenes, and world-coordinate global human trajectory reconstruction.
\end{enumerate}


% ===================================================================
% 2. RELATED WORK
% ===================================================================
\section{Related Work}
\label{sec:related}

\subsection{Monocular 3D Human Pose and Shape Estimation}
Methods for recovering 3D human pose and shape from a single image fall into optimization-based and regression-based categories. Optimization methods~\cite{SMPLify2016,NIKI2023} fit SMPL~\cite{SMPL} parameters by minimizing 2D keypoint reprojection errors but incur high computational cost and initialization sensitivity. Regression methods, pioneered by HMR~\cite{HMR2018}, directly predict SMPL parameters from image features. Subsequent works~\cite{PyMAF2021,CLIFF2022,TokenHMR2024} have steadily improved per-frame accuracy through better architectures and training strategies. Recent foundation-model-scale approaches~\cite{HMR2024,SMPLerX2024,Sapiens2024} demonstrate that scaling up model capacity and training data yields further gains, and diffusion-based methods~\cite{ScoreHMR2024} leverage learned priors for test-time refinement. Multi-person extensions~\cite{MultiHMR2024} recover all humans in a single forward pass.

When applied to video, per-frame methods lack temporal consistency. Video-based approaches~\cite{TCMR2021,MPSNet2022, MotionBERT2023,kanazawa2019learning, luo20203d} integrate multi-frame information through temporal architectures to enhance smoothness, yet often sacrifice per-frame accuracy---revealing an inherent tension between temporal coherence and pixel-level precision. More critically, all these methods operate in the camera coordinate frame, making it difficult to recover physically plausible global motion under moving cameras.

\subsection{Global Human Motion Recovery from Moving Cameras}

Recovering global human trajectories from moving-camera videos requires jointly reasoning about camera and human motion. GLAMR~\cite{GLAMR2022} predicts global trajectories through learned motion priors but ignores camera cues, suffering from trajectory drift. SLAHMR~\cite{SLAHMR2023} and PACE~\cite{PACE2024} introduce SLAM for camera estimation and perform joint optimization by minimizing the negative log-likelihood of motion priors, but depend on MoCap-trained distributions~\cite{AMASS2019} and incur high computational cost. HuMoR~\cite{HuMoR2022} provides a generative motion prior for robust pose estimation but does not model camera motion. WHAM~\cite{WHAM2024} achieves efficient estimation via IMU feature fusion but accumulates angular velocity errors. GVHMR~\cite{GVHMR2024} predicts motion in gravity-view coordinates to reduce drift but still relies on a separate visual odometry pass. TRACE~\cite{TRACE2023} adopts end-to-end regression with high efficiency but weak temporal consistency. TRAM~\cite{TRAM2024} proposes a two-stage pipeline using masked DROID-SLAM~\cite{DROIDSLAM2021} for camera estimation, but masking out human regions severely degrades SLAM accuracy in human-dominated scenes. Diffusion-based approaches~\cite{RoHM2024,PhysDiff2023} refine motion through learned priors but do not explicitly model the camera-human coupling. Concurrent work HumanMM~\cite{HumanMM2025} addresses multi-shot videos but assumes shot boundaries are given.

A common limitation of these methods is the unidirectional information flow between camera and human estimation. In contrast, EIHMR reverses the role of human motion in the pipeline: rather than treating dynamic humans as obstacles to be masked out, it renders refined human motion back into the images to provide structured geometric cues for SLAM, transforming dynamic humans from a source of interference into a source of useful constraints.


\begin{figure*}[!t]
  \centering
  \includegraphics[width=\linewidth]{images/Network.pdf}
  \caption{\textbf{Overview of EIHMR.} Given a monocular moving-camera video, we first obtain coarse camera parameters and initial SMPL parameters. \textbf{Stage~1 (SA-LHMR):} For each keyframe, neighboring frames within a sliding window are projected into the keyframe's viewpoint. A two-stage refinement corrects the root trajectory and local motion sequences. \textbf{Stage~2 (MO-SLAM):} The refined local meshes are rendered back onto the original keyframe images with golden-angle HSV coloring, producing motion-augmented images. The motion-augmented images are fed into the motion-aware SLAM method to recover refined camera trajectories. The final global human motion is obtained by combining the refined cameras with SMPL parameters.}
  \label{fig:overview}
\end{figure*}

% ===================================================================
% 3. METHOD
% ===================================================================
% ===================================================================
% 3. METHOD
% ===================================================================
\section{Method}
\label{sec:method}

Our goal is to recover complete 3D human motion from monocular video $\mathcal{V} = \{I_t\}_{t=0}^{T}$ captured by an in-the-wild moving camera. We decompose the global motion into root trajectory $\{Trj_t\}_{t=0}^{T}$ in the world coordinate system and a kinematic body motion represented by SMPL~\cite{SMPL,SMPLX2019} pose sequences $\{\Theta_t\}_{t=0}^{T}$ in the camera coordinate system. The global human motion can be obtained by estimating the camera trajectory $\{C_t\}_{t=0}^{T}$ and the human position relative to the camera at each time step $\{H_t\}_{t=0}^{T}$:
%
\begin{equation}
  Trj_t = C_t \cdot H_t.
\end{equation}

\subsection{EIHMR Pipeline}
\label{sec:pipeline}

The prevailing two-stage paradigm first estimates camera trajectories by masking out dynamic human regions to satisfy the static-scene assumption of SLAM:
%
\begin{equation}
  \{C_t\} = \text{SLAM}(\{I_t\}, \{M_t\}),
\end{equation}
%
and then reconstructs human body motion conditioned on the recovered camera parameters with HMR method:
%
\begin{equation}
  \{\Theta_t, H_t\} = \text{HMR}(\mathcal{V}, \{C_t\}).
\end{equation}
%
This unidirectional pipeline treats camera estimation and motion reconstruction as isolated processes: SLAM receives no benefit from the estimated human motion, while the human pose estimator inherits all camera errors without any mechanism for correction.

EIHMR replaces this isolation with a decouple-then-collaborate loop (\cref{fig:overview}). SA-LHMR decouples camera and human motion by projecting pose sequences into frozen keyframe viewpoints, where kinematic constraints and metric depth are applied for scene-aware local refinement:
%
\begin{equation}
  \{\Theta_t', H_t'\} = \text{SA-LHMR}\left(\mathcal{V}, \{C_t\}, \{\Theta_t\}, \{H_t\}\right).
\end{equation}

The refined motion is subsequently re-rendered as geometrically consistent static meshes onto the original frames, producing motion-augmented images $\{\tilde{I}_t\}$. MO-SLAM leverages these structured visual cues for robust camera estimation:
%
\begin{equation}
  \{C_t'\} = \text{MO-SLAM}\left(\{\tilde{I}_t\}, \{M_t\}\right).
\end{equation}

The detailed formulations of SA-LHMR and MO-SLAM are presented in \cref{sec:sa-lhmr} and \cref{sec:mo-slam}, respectively.



\subsection{Scene-Aware Local Human Motion Reconstruction (SA-LHMR)}
\label{sec:sa-lhmr}

We select a set of keyframes $\mathcal{K} = \{k_i\}$ from the video and use each keyframe's viewpoint as a local reference frame. For keyframe $k_i$, the human motion of neighboring frames within a sliding window $\mathcal{W}_i$ is projected into this viewpoint via camera parameters $\{C_t\}$, yielding local trajectory sequences $\{H_{ij}^{k}\}$ and body poses $\{\Theta_{ij}\}$, where $i$ is the keyframe index and $j$ is the frame index within the window. Since all frames now share a unified viewpoint, camera ego-motion is completely eliminated, and many non-physical artifacts in the motion estimation become clearly discernible: the human model lacks proper contact with the scene ground, foot joints that should remain stationary exhibit kinematically implausible skating displacements, and limbs penetrate scene surfaces such as stairs and floors. These non-physical motion artifacts arise from two sources: first, estimation errors in the initial SLAM-provided camera parameters; second, the temporally-aware HPE network sacrificing geometric consistency with image observations to maintain temporal smoothness under erroneous camera conditioning. Notably, these errors can be explicitly measured and corrected through forward kinematics and scene depth, yet are difficult for end-to-end deep networks to learn effectively.

We propose a two-stage motion refinement strategy. Firstly, we correct the root trajectory of the motion sequence based on static joint analysis from forward kinematics. Then, we detect scene-penetrating contact points using metric depth maps and optimize body poses via inverse kinematics to eliminate penetrations. Crucially, since all frames under the unified viewpoint share a single depth map reference, we can constrain multiple contact points across multiple frames within the entire window using a single depth estimate, forming a powerful spatiotemporal consistency constraint.

\paragraph{Root Trajectory Correction via Forward Kinematics.}
We traverse the local sequence under keyframe $k_i$ in a sliding-window manner. For a window $\mathcal{W}_i = [k_i{-}W, k_i{+}W]$ of size $2W{+}1$ anchored at center frame $k_i$, we first compute the keyframe-coordinate joint positions for each frame within the window via forward kinematics (FK):
%
\begin{equation}
  J_{ij} = \text{FK}(\Theta_{ij}, H_{ij}^{k}), \quad j \in \mathcal{W}_i.
\end{equation}

The HPE network simultaneously outputs per-frame static confidence scores $\mathbf{s}_j \in [0,1]^6$, which are used to identify joints that should maintain static contact with the scene. For joints identified as static, we compute their inter-frame displacements and use them as correction signals for the root trajectory: since static joints should not undergo displacement in the world coordinate system, the inverse of their displacement is back-propagated to the root trajectory $H_{ij}^{k}$ to eliminate trajectory deviations introduced by camera estimation errors and HPE temporal smoothing. The correction process is anchored at the window center frame, keeping its root position unchanged and only adjusting the relative displacements of neighboring frames.

For frames in overlapping window regions, the corrected root trajectories contributed by each window are merged via normalized Gaussian weights based on the frame-to-center distance:
%
\begin{equation}
  H_{j}^{k'} = \frac{\sum_{i} w_{ij} \cdot H_{ij}^{k'}}{\sum_{i} w_{ij}}, \quad w_{ij} = \exp\!\left(-\frac{(j - c_i)^2}{2\sigma^2}\right),
\end{equation}
%
where $c_i$ is the center frame index of the $i$-th window and $\sigma$ controls the Gaussian decay bandwidth, ensuring smooth transitions across overlapping regions.

\paragraph{Contact-Aware Inverse Kinematics via Scene Depth.}
After root trajectory correction, the motion sequence may still contain non-physical phenomena where limbs penetrate scene surfaces---for example, feet passing through stair steps or the ground plane. We leverage the metric depth map $D_c = \text{Metric3D}(I_{k_i})$ estimated from the keyframe viewpoint to detect and correct these penetrations.

First, we perform window-local scale alignment. Since scale discrepancies may exist between the SLAM-estimated camera coordinate system and the metric depth map, we project the 3D joints $J_{ij} \in \mathbb{R}^3$ in the keyframe camera coordinate system onto the image plane to obtain its pixel coordinates and depth:
%
\begin{equation}
  z_j^{\text{smpl}} \begin{pmatrix} u_j \\ v_j \\ 1 \end{pmatrix} = K \cdot J_{ij},
\end{equation}
%
where $K$ is the camera intrinsic matrix, $(u_j, v_j)$ are the projected pixel coordinates, and $z_j^{\text{smpl}}$ is the joint depth. The corresponding scene depth $z_j^{\text{scene}} = D_c(u_j, v_j)$ is sampled from the metric depth map at $(u_j, v_j)$. Using all 22 joints of the center frame as sample points, we compute a scale factor:
%
\begin{equation}
  \alpha = \operatorname{median}\!\left(\frac{z_j^{\text{scene}}}{z_j^{\text{smpl}}}\right), \quad j = 1, \ldots, 22.
\end{equation}

Subsequently, for each static foot joint in every frame of the window, we check whether it penetrates the scene surface: the joint is projected into the keyframe's camera coordinate system, and its scale-aligned depth $z_j^{\text{smpl}} \cdot \alpha$ is compared against the depth map value $z_j^{\text{scene}}$. If $z_j^{\text{smpl}} \cdot \alpha > z_j^{\text{scene}} + \epsilon$, where $\epsilon$ is the penetration tolerance, the joint is deemed to be penetrating. For each penetrating joint, we obtain its target position in world coordinates via depth map back-projection:
%
\begin{equation}
  \mathbf{p}_j^{\text{target}} = \frac{z_j^{\text{scene}}}{\alpha} \cdot K^{-1} \begin{pmatrix} u_j \\ v_j \\ 1 \end{pmatrix}, 
\end{equation}
%
where $(u_j, v_j)$ are the joint's pixel coordinates. Finally, we optimize the pose increments $\delta\boldsymbol{\theta}$ along the leg kinematic chains (hip $\to$ knee $\to$ ankle $\to$ toe) for all frames within the window, minimizing the distance from contact joints to their target positions:
%
\begin{equation}
  \mathcal{L}_{\text{contact}} = \sum_{j \in \mathcal{S}} \rho\!\left(\left\|\text{FK}(\Theta_j + \delta\boldsymbol{\theta})_j - \mathbf{p}_j^{\text{target}}\right\|_2\right), 
\end{equation}
%
where $\mathcal{S}$ is the set of all penetrating static joints and $\rho(\cdot)$ is the Geman--McClure robust kernel. The optimization simultaneously enforces pose regularization $\mathcal{L}_{\text{reg}} = \|\delta\boldsymbol{\theta}\|_2^2$ and a temporal smoothness constraint $\mathcal{L}_{\text{smooth}}$ to ensure that the corrected poses are physically plausible and temporally coherent. Analogous to the root trajectory correction, the optimized body poses $\Theta_{ij}'$ from each window are merged via Gaussian distance weighting in overlapping regions to produce a globally consistent pose sequence.



\subsection{Motion-Aware Global SLAM (MO-SLAM)}
\label{sec:mo-slam}

With geometrically consistent local motion from SA-LHMR, we now leverage it to improve camera trajectory estimation. Traditional SLAM systems~\cite{ORBSLAM3,DROIDSLAM2021,DPVO2023} assume a static scene, requiring dynamic objects to be masked or down-weighted and thereby discarding potentially useful information. Recent feed-forward 3D reconstruction methods~\cite{DUSt3R2024,MonST3R2024,MASt3RSLAM2025} handle dynamic scenes implicitly but do not model human motion structure. As shown in Stage 2 of \cref{fig:overview}, the core idea of MO-SLAM is to render the refined motion back into the video frames so that dynamic human regions, which normally degrade SLAM, instead provide structured geometric cues for inter-frame feature matching.

To construct the motion-augmented images, we render the refined motion from SA-LHMR back into each keyframe's image. For keyframe $k_i$, we retrieve its local root trajectories $\{H_{ij}^{k'}\}$ and the globally merged body poses $\{\Theta_j'\}$ for all frames $j$ within the window $\mathcal{W}_i$. Each SMPL mesh is projected onto the image plane of $k_i$ via the keyframe's camera intrinsics and rendered as a solid overlay. To provide SLAM with discriminative cross-frame correspondences, we assign each frame $j$ a distinct color using the golden-angle HSV scheme:
%
\begin{equation}
  \text{color}(j) = \text{HSV}\!\left(\frac{j \cdot \phi}{2\pi} \bmod 1\right), \quad \phi = \frac{2\pi}{\varphi^2},
\end{equation}
%
where $\varphi = \frac{1+\sqrt{5}}{2}$ is the golden ratio. This scheme maximizes the perceptual distance between colors of consecutive frames, ensuring that SLAM can establish unambiguous inter-frame feature correspondences from the rendered meshes. In the resulting motion-augmented images $\{\tilde{I}_t\}$, the originally dynamic human regions are transformed into multi-view geometrically consistent static structures with frame-distinguishing color codes.

We then feed these motion-augmented images into DROID-SLAM \cite{DROIDSLAM2021} with three targeted modifications. First, rendered frame indices are injected as forced keyframes by modifying the MotionFilter to bypass its optical-flow-based~\cite{RAFT2020} keyframe selection criterion and directly admit frames from the forced set; the Frontend is further prevented from removing any keyframes, ensuring that all motion-augmented frames participate in bundle adjustment. Second, we introduce an adaptive image replacement strategy: at each keyframe we compute the grayscale Laplacian variance as a texture score, and when this score falls below a threshold $\tau$ the rendered image replaces the original to provide additional structural constraints, while the original masked image is retained otherwise, allowing MO-SLAM to adapt to varying scene complexity. Third, even with rendered images, we apply human-aware correlation weighting by using human masks to down-weight human regions in the correlation computation of SLAM; although the rendered SMPL meshes provide correct geometric contours, their surface appearance differs from the real image, making high pixel-level matching weights inappropriate. This design exploits motion information at the geometry level while suppressing it at the pixel level.


% ===================================================================
% 4. EXPERIMENTS
% ===================================================================
\section{Experiments}
\label{sec:experiments}

\subsection{Setup}
\label{sec:impl}

\noindent\textbf{Implementation Details.}
We use PromptHMR~\cite{PromptHMR2024} as the pose estimation backbone, with YOLOv7~\cite{YOLO} for detection. SA-LHMR uses half-window size $W{=}10$, corresponding to a window size of $2W{+}1{=}21$, stride $s{=}5$, and Gaussian $\sigma{=}5$, optimizing each window for 50 iterations with Adam. Metric depth is from Metric3D~\cite{Metric3D2024} computed every 5 frames and 2D keypoints from ViTPose~\cite{ViTPose2022}. MO-SLAM is built on DROID-SLAM~\cite{DROIDSLAM2021} with rendering interval 5. Human masks are provided by SAM~\cite{SAM2023}.

\noindent\textbf{Datasets.}
\textbf{EMDB}~\cite{EMDB2023} provides EMDB1 for simple scenarios and EMDB2 containing 25 egocentric head-mounted videos with camera poses and SMPL ground truth. EMDB2 is the primary testbed; following TRAM~\cite{TRAM2024}, we split sequences into Short/Medium/Long by camera travel distance.
\textbf{3DPW}~\cite{3DPW2018} is an outdoor dataset with IMU-assisted SMPL annotations for generalization evaluation.

\noindent\textbf{Metrics.}
\emph{Camera}: Absolute Trajectory Error, ATE in meters, after rigid alignment with scaling.
\emph{Body pose}: PA-MPJPE, MPJPE, PVE, all in mm, and Accel in m/s$^2$.
\emph{Global trajectory}: W-MPJPE in mm aligned on the first two frames, WA-MPJPE in mm with full-segment alignment, and RTE in \% as root translation error normalized by total displacement, evaluated on 100-frame segments. Body pose is reported on EMDB and 3DPW; global trajectory and camera metrics on EMDB2 only.


% ------------------------------------------------------------------
% 4.2 Body Pose Evaluation
% ------------------------------------------------------------------
\subsection{Body Pose Evaluation}
\label{sec:exp-pose}

\begin{table}[t]
  \caption{\textbf{Body pose evaluation on EMDB and 3DPW.} Parentheses denote the number of body joints. Accel in m/s$^2$, others in mm. Best in \textbf{bold}, second best \underline{underlined}.}
  \label{tab:pose}
  \centering
  \resizebox{\columnwidth}{!}{%
  \begin{tabular}{@{}lcccccccc@{}}
    \toprule
    & \multicolumn{4}{c}{3DPW (14)} & \multicolumn{4}{c}{EMDB (24)} \\
    \cmidrule(lr){2-5} \cmidrule(lr){6-9}
    Method & PA-MPJPE$\downarrow$ & MPJPE$\downarrow$ & PVE$\downarrow$ & Accel$\downarrow$ & PA-MPJPE$\downarrow$ & MPJPE$\downarrow$ & PVE$\downarrow$ & Accel$\downarrow$ \\
    \midrule
    TCMR~\cite{TCMR2021} & 52.7 & 86.5 & 101.4 & 6.0 & 79.8 & 127.7 & 150.2 & \underline{5.3} \\
    VIBE~\cite{VIBE2020} & 51.9 & 82.9 & 98.4 & 18.5 & 81.6 & 126.1 & 149.9 & 26.5 \\
    MPS-Net~\cite{MPSNet2022} & 52.1 & 84.3 & 99.0 & 6.5 & 81.4 & 123.3 & 143.9 & 6.2 \\
    GLoT~\cite{GLoT2023} & 50.6 & 80.7 & 96.4 & 6.0 & 79.1 & 119.9 & 140.8 & 5.4 \\
    GLAMR~\cite{GLAMR2022} & 51.1 & -- & -- & 8.0 & 73.8 & 113.8 & 134.9 & 33.0 \\
    TRACE~\cite{TRACE2023} & 50.9 & 79.1 & 95.4 & 28.6 & 71.5 & 110.0 & 129.6 & 25.5 \\
    WHAM~\cite{WHAM2024} & 37.5 & 59.8 & 71.5 & 6.6 & 50.4 & 79.7 & 94.4 & \underline{5.3} \\
    TRAM~\cite{TRAM2024} & \underline{36.7} & 59.3 & \underline{69.6} & \textbf{4.9} & 45.7 & 74.4 & 86.6 & \textbf{4.9} \\
    PromptHMR~\cite{PromptHMR2024} & \textbf{36.6} & \underline{58.7} & 69.7 & 24.8 & \underline{43.1} & \textbf{69.6} & \underline{82.8} & 20.5 \\
    \textbf{EIHMR (Ours)} & 36.9 & \textbf{57.9} & \textbf{69.4} & \underline{5.5} & \textbf{42.6} & \underline{72.3} & \textbf{82.4} & 8.3 \\
    \bottomrule
  \end{tabular}%
  }
\end{table}

\noindent\textbf{Quantitative results.}
As shown in \cref{tab:pose}, EIHMR's primary gain over PromptHMR is temporal stability: SA-LHMR reduces Accel by ${\sim}$60\%--78\% on both benchmarks, reaching parity with dedicated video methods. EIHMR also achieves the best PVE on both datasets and the best PA-MPJPE of 42.6\,mm on EMDB. PromptHMR leads on EMDB MPJPE (69.6 vs.\ 72.3\,mm); this trade-off is expected, as contact-aware IK shifts foot joints toward the scene surface, slightly perturbing Procrustes-optimal alignment but yielding physically grounded poses that translate into a 10\% W-MPJPE gain in \cref{tab:global}. EIHMR's Accel of 8.3 on EMDB remains above TRAM's 4.9 due to TRAM's dedicated temporal architecture, yet the ${\sim}$60\% reduction over PromptHMR confirms that physics-based constraints effectively bridge the single-frame-to-video gap.

\begin{figure*}[!t]
  \centering
  \includegraphics[width=\linewidth]{images/knee_pose}
  \caption{\textbf{Qualitative body pose comparison on stair walking.} Top: input frames. Bottom: left knee flexion curve (frames 815--885); GT (\textcolor{gray}{gray}), Ours (\textcolor{green!50!black}{green}), PromptHMR (\textcolor{blue}{blue}), TRAM (\textcolor{orange}{orange}). Dashed lines mark two flexion peaks (frames 823 and 858), with corresponding SMPL leg meshes shown alongside each peak.}
  \label{fig:qualitative}
\end{figure*}

\noindent\textbf{Qualitative results.}
\cref{fig:qualitative} visualizes left knee flexion on a stair-descending sequence. The flexion curve reveals that at both peaks, frames 823 and 858, EIHMR closely tracks the ground truth, whereas PromptHMR~\cite{PromptHMR2024} and TRAM~\cite{TRAM2024} consistently overestimate the flexion angle by a clear margin. The accompanying SMPL meshes provide further insight. At frame 823, EIHMR produces a knee configuration that best matches GT; PromptHMR over-flexes the knee, yielding an incorrect joint angle, and TRAM's peak has not yet arrived due to temporal lag. At frame 858, EIHMR and PromptHMR show similar overall leg poses, but EIHMR achieves more accurate foot placement, consistent with SA-LHMR's contact-aware IK pulling foot joints onto the scene surface, while TRAM again exhibits a delayed response, with its flexion peak shifted later than GT. These results demonstrate that SA-LHMR's scene-aware optimization corrects both joint-angle accuracy at contact events and temporal alignment, while reducing Accel from 20.5 to 8.3\,m/s$^2$ over the PromptHMR backbone.


% ------------------------------------------------------------------
% 4.3 Camera Trajectory Evaluation
% ------------------------------------------------------------------
\subsection{Camera Trajectory Evaluation}
\label{sec:exp-camera}

All camera evaluations are conducted on EMDB2, where humans occupy a large portion of the frame and severely challenge standard SLAM systems.
\begin{figure*}[!t]
  \centering
  \includegraphics[width=\linewidth]{images/19_camera_pose_tra}
  \caption{\textbf{Camera trajectory evaluation on EMDB2.} Left: ATE (m) across Short/Medium/Long/Avg splits. Right: 3D camera trajectory (top) and global body trajectory (bottom) on a representative sequence, with inset frames showing the corresponding video input. GT (\textcolor{gray}{gray dashed}), EIHMR (\textcolor{green!50!black}{green}), PromptHMR (\textcolor{cyan}{blue}). Triangle and circle markers denote trajectory start and end points, respectively.}
  \label{fig:cam-body-traj}
\end{figure*}

\noindent\textbf{Quantitative results.}
As shown in \cref{fig:cam-body-traj} (left), DROID-SLAM~\cite{DROIDSLAM2021} and ORB-SLAM~\cite{ORBSLAM3} suffer severe degradation on longer sequences, with ATE exceeding 2\,m and 3\,m on Long splits respectively, as neither handles dynamic human regions. EIHMR and TRAM~\cite{TRAM2024} both reduce ATE below 0.5\,m across all splits. EIHMR achieves the lowest Short ATE; TRAM holds a slight edge on average ATE because its masking strategy completely removes human pixels that inevitably carry some residual appearance mismatch between the rendered SMPL surface texture and the real image. Importantly, camera ATE is measured after rigid alignment with scaling, which absorbs global scale and orientation offsets but does not capture the local geometric consistency between camera and human motion that ultimately determines global trajectory quality. This distinction explains why EIHMR's marginally higher average ATE nonetheless yields substantially better global trajectory metrics in \cref{sec:exp-global}: MO-SLAM preserves human-region geometric structure during bundle adjustment, producing camera trajectories that are metrically consistent with the human body scale and position, a property not captured by ATE's rigid alignment but directly reflected in W-MPJPE and WA-MPJPE.

\noindent\textbf{Qualitative results.}
\cref{fig:cam-body-traj} (top-right) shows that EIHMR's camera trajectory closely follows the ground truth, while PromptHMR~\cite{PromptHMR2024} accumulates progressive drift without explicit human motion modeling. The bottom-right panel further illustrates the resulting global body trajectory: EIHMR maintains close alignment with GT throughout the sequence, whereas PromptHMR drifts substantially. This confirms that rendering refined motion into frames stabilizes SLAM in human-dominated egocentric scenes, and that accurate camera trajectories from MO-SLAM directly yield faithful global body trajectories.



% ------------------------------------------------------------------
% 4.4 Global Human Trajectory Evaluation
% ------------------------------------------------------------------
\subsection{Global Human Trajectory Evaluation}
\label{sec:exp-global}

Global trajectory recovery combines camera trajectories from MO-SLAM with local poses from SA-LHMR by transforming camera-frame root positions into world coordinates.

\begin{table}[t]
  \caption{\textbf{Global human trajectory evaluation on EMDB2.} PA-MPJPE, WA-MPJPE, W-MPJPE in mm; RTE in \%. Best in \textbf{bold}, second best \underline{underlined}.}
  \label{tab:global}
  \centering
  \setlength{\tabcolsep}{6pt}
  \begin{tabular}{@{}lcccc@{}}
    \toprule
    Method & PA-MPJPE$\downarrow$ & WA-MPJPE$\downarrow$ & W-MPJPE$\downarrow$ & RTE$\downarrow$ \\
    \midrule
    GLAMR~\cite{GLAMR2022} & 56.00 & 280.80 & 726.60 & 11.40 \\
    SLAHMR~\cite{SLAHMR2023} & 61.50 & 326.90 & 776.10 & 10.20 \\
    WHAM~\cite{WHAM2024} & 38.20 & 133.30 & 343.90 & 4.60 \\
    TRAM~\cite{TRAM2024} & 38.21 & \underline{85.35} & \underline{263.35} & \textbf{1.58} \\
    PromptHMR~\cite{PromptHMR2024} & \textbf{34.48} & 88.34 & 258.10 & 2.62 \\
    \textbf{EIHMR (Ours)} & \underline{35.55} & \textbf{81.34} & \textbf{232.22} & \underline{2.56} \\
    \bottomrule
  \end{tabular}
\end{table}

\noindent\textbf{Quantitative results.}
As shown in \cref{tab:global}, EIHMR achieves the best WA-MPJPE of 81.34\,mm and W-MPJPE of 232.22\,mm, reducing global trajectory error by ${\sim}$10\% over PromptHMR~\cite{PromptHMR2024} and ${\sim}$11.8\% over TRAM~\cite{TRAM2024}, while substantially outperforming GLAMR, SLAHMR, and WHAM. PromptHMR achieves the best PA-MPJPE at 34.48 vs.\ 35.55\,mm, consistent with its higher per-frame accuracy on EMDB in \cref{tab:pose}; however, PA-MPJPE applies per-frame Procrustes alignment that factors out global translation, and thus does not reflect trajectory quality. PromptHMR's W-MPJPE is 11\% worse at 258 vs.\ 232\,mm, confirming that strong local accuracy alone cannot compensate for camera drift. TRAM achieves the best RTE of 1.58\% owing to its low average camera ATE; RTE normalizes root translation error by total displacement and is thus dominated by camera trajectory accuracy, where TRAM's mask strategy excels as discussed in \cref{sec:exp-camera}. By contrast, EIHMR optimizes for geometric consistency between the human body and the scene, which is better captured by W-MPJPE and WA-MPJPE, the two metrics that jointly evaluate the accuracy of full 3D joint positions in world coordinates. Overall, EIHMR surpasses TRAM despite the latter employing masked DROID-SLAM, demonstrating that rendering motion back into frames is more effective than discarding human regions entirely for recovering geometrically faithful global trajectories.

\begin{figure*}[!t]
  \centering
  \includegraphics[width=\linewidth]{Qualitative}
  \caption{\textbf{Qualitative global motion comparison.} GT (\textcolor{gray}{gray}), EIHMR (\textcolor{green!50!black}{green}), PromptHMR (\textcolor{blue}{blue}) on two EMDB sequences. Top row: input frames with SMPL overlays; bottom row: 3D trajectories from two viewpoints. Red-box zoom-ins (bottom-right) highlight foot-level details: GT and EIHMR exhibit coherent, grounded footsteps, whereas PromptHMR shows visible foot sliding and floating artifacts.}
  \label{fig:global-qual}
\end{figure*}

\begin{table}[t]
  \caption{\textbf{Ablation study on EMDB2.} Each row removes one component. PA-MPJPE, WA-MPJPE, W-MPJPE in mm; RTE in \%; ATE in m. All $\downarrow$.}
  \label{tab:ablation}
  \centering
  \setlength{\tabcolsep}{4pt}
  \resizebox{\columnwidth}{!}{%
  \begin{tabular}{@{}lccccc@{}}
    \toprule
    Configuration & PA-MPJPE$\downarrow$ & WA-MPJPE$\downarrow$ & W-MPJPE$\downarrow$ & RTE$\downarrow$ & ATE$\downarrow$ \\
    \midrule
    \textbf{Full model (EIHMR)} & 35.55 & \textbf{81.34} & \textbf{232.22} & 2.56 & \textbf{0.42} \\
    \quad $-$ Contact-Aware IK & \textbf{35.43} & 84.87 & 245.23 & 2.66 & 0.42 \\
    \quad $-$ FK Root Correction & 40.56 & 89.52 & 266.52 & 3.12 & 0.65 \\
    \quad $-$ Adaptive Image Replacement & 35.56 & 102.58 & 278.67 & \textbf{2.37} & 0.57 \\
    \quad $-$ Timestep Color Encoding & 35.56 & 129.47 & 342.21 & 3.13 & 0.84 \\
    \bottomrule
  \end{tabular}%
  }
\end{table}


\noindent\textbf{Qualitative results.}
\cref{fig:global-qual} visualizes global motion on two EMDB sequences. PromptHMR~\cite{PromptHMR2024} produces progressively drifted trajectories, showing that even a strong pose backbone with PA-MPJPE 43.1 on EMDB as reported in \cref{tab:pose} cannot compensate for inaccurate camera trajectories. EIHMR maintains close alignment with the ground truth across both viewpoints. The red-box zoom-ins further reveal foot-level quality: GT and EIHMR produce coherent, grounded footsteps, while PromptHMR exhibits clear foot sliding and floating, confirming that SA-LHMR's contact-aware constraints translate into physically plausible global motion. \cref{fig:cam-body-traj} further confirms this causal chain: accurate camera trajectories from MO-SLAM directly yield faithful global body trajectories.


% ------------------------------------------------------------------
% 4.5 Ablation Study
% ------------------------------------------------------------------
\subsection{Ablation Study}
\label{sec:ablation}

We ablate each component on EMDB2, as shown in \cref{tab:ablation}.


\noindent\textbf{SA-LHMR components.}
Removing Contact-Aware IK marginally improves PA-MPJPE (35.55$\to$35.43) as the IK adjustment can shift joints from Procrustes-optimal alignment, but degrades W-MPJPE by 5.6\%, confirming its primary role in enforcing foot-ground contact for global recovery. Removing FK Root Correction causes the largest degradation: PA-MPJPE rises 14\%, ATE rises 55\%, and W-MPJPE rises 15\%, confirming it as the foundational anchor bridging local pose quality and global trajectory accuracy.

\noindent\textbf{MO-SLAM components.}
Removing Timestep Color Encoding produces the most severe camera degradation: ATE doubles to 0.84\,m, W-MPJPE rises 47\%, and WA-MPJPE rises 59\%, confirming that temporal distinctiveness in rendered cues is essential for SLAM. Removing Adaptive Image Replacement degrades ATE by 36\% and WA-MPJPE by 26\% while leaving PA-MPJPE unchanged. Notably, RTE slightly improves from 2.56 to 2.37\%: since RTE normalizes by total displacement, the replacement strategy may marginally bias SLAM scale estimation; however, this minor regression is far outweighed by the 26\% WA-MPJPE improvement.


% ===================================================================
% 5. CONCLUSION
% ===================================================================
\section{Conclusion}
\label{sec:conclusion}


We presented EIHMR, a collaborative human--camera co-estimation framework that replaces the isolated, unidirectional pipeline of prior work with a decouple-then-collaborate loop. SA-LHMR decouples camera and human motion via frozen-keyframe reprojection and enforces scene consistency through forward-kinematics root correction and depth-guided contact-aware inverse kinematics, eliminating foot sliding and scene penetration. MO-SLAM re-renders the refined motion back into the original frames, transforming dynamic human regions from masked interference into structured geometric cues for camera estimation. This bidirectional information flow allows human motion and camera trajectory to mutually reinforce each other. Experiments on EMDB show that EIHMR reduces global trajectory error by approximately 10\% over the strongest baselines, with simultaneous gains in both body pose and camera estimation, confirming that leveraging human motion as geometric constraints is more effective than discarding it.

\noindent\textbf{Limitations.}
Performance may degrade under extreme camera motions (\eg, severe pitch, large rotations) or when depth cues are insufficient due to textureless clothing. Future work could incorporate temporal motion priors~\cite{MDM2023,PhysDiff2023} or multi-modal inputs (\eg, IMU~\cite{WHAM2024}) to improve robustness, and extend the framework to multi-person scenarios~\cite{MultiHMR2024}.


% ===================================================================
% ACKNOWLEDGEMENTS & BIBLIOGRAPHY
% ===================================================================

%\section*{Acknowledgements}
%TODO: Add acknowledgements here for the camera-ready version.

\bibliographystyle{splncs04}
\bibliography{main}

\end{document}
